\documentclass[11pt]{article}
\usepackage[a4paper,margin=1in]{geometry}
\usepackage{amsmath,amssymb,mathtools,bm}
\usepackage{enumitem,booktabs,array}
\usepackage{tcolorbox}
\usepackage{hyperref}
\usepackage[protrusion=true,expansion=false]{microtype}
\hypersetup{colorlinks=true,linkcolor=blue,urlcolor=blue,citecolor=blue}
\setlist[itemize]{topsep=2pt,itemsep=2pt}
\setlist[enumerate]{topsep=2pt,itemsep=2pt}
\newtcolorbox{keybox}{colback=blue!3!white,colframe=blue!40!black,boxrule=0.4pt,arc=1mm}
\newtcolorbox{alertbox}{colback=red!3!white,colframe=red!40!black,boxrule=0.4pt,arc=1mm}
\newtcolorbox{answerbox}{colback=green!3!white,colframe=green!40!black,boxrule=0.4pt,arc=1mm}
\newtcolorbox{drillbox}{colback=yellow!5!white,colframe=orange!60!black,boxrule=0.4pt,arc=1mm}
\newcommand{\EV}{\mathbb{E}}
\newcommand{\Var}{\mathrm{Var}}
\newcommand{\Cov}{\mathrm{Cov}}
\newcommand{\blank}[1]{\underline{\hspace{#1}}}

\title{W14-01: Li, Shan, Tang \& Yao (2024, RFS) \\
\large ``Corporate Climate Risk: Measurements and Responses'' --- Active Recall Workbook}
\author{Banking \& Risk Management --- Exam Prep}
\date{\today}
\begin{document}\maketitle

\begin{keybox}
\textbf{Paper in one sentence.} LSTY build firm-level climate-risk measures (physical, transition, regulatory) from 10-K text using supervised ML trained on a dictionary of climate terms, and show firms with higher measured climate risk \emph{raise cash} holdings, \emph{cut capex}, and \emph{alter capital structure} --- real adjustments consistent with precautionary-savings and optionality-value arguments.

\textbf{Placement.} Text-measurement-of-climate-risk paper; extends Sautner et al.\ 2023 and related climate-risk-text work. Pairs with KSS (W13-03) on institutional climate attitudes.
\end{keybox}

\section*{Round 1: Complete Walk-Through}

\subsection*{1.1 Measurement}
Parse 10-K text; classify climate-related sentences into physical, transition, and regulatory risks. Supervised ML with hand-coded training set. Construct firm-year climate-risk score for 2000--2020. Sample $\sim$6{,}000 US firms.

\subsection*{1.2 Validation}
\begin{itemize}
\item Correlates with objective exposures (coastal siting, emissions-intensive industries).
\item Correlates with analyst-discussed climate sections in conference calls.
\item Correlates with CDP/TCFD disclosure intensities.
\end{itemize}

\subsection*{1.3 Central tests}
Panel regression with firm and year FE:
\[
Y_{it}=\alpha_i+\delta_t+\beta\cdot\text{ClimateRisk}_{it-1}+X_{it-1}'\gamma+\varepsilon_{it}.
\]
$Y_{it}$ includes cash/assets, capex/assets, leverage, dividend policy.

\subsection*{1.4 Headline results}
\begin{itemize}
\item Firms with higher lagged climate-risk score hold more cash (+0.5--1 pp of assets).
\item They cut capex by $\sim$0.3--0.5 pp.
\item Leverage declines modestly.
\item Dividend payout is less affected.
\item Effects are concentrated in physical-risk-exposed firms.
\end{itemize}

\subsection*{1.5 Interpretation}
\begin{itemize}
\item Firms \emph{actively manage} perceived climate risk via financial flexibility (BKS 2009 precautionary cash extended to climate).
\item Investment option-value: wait-and-see under regulatory uncertainty reduces capex.
\item Complements BMFS (2016) on hedging adjustments and FRS (2021) on flexibility premium.
\end{itemize}

\subsection*{1.6 Caveats}
\begin{alertbox}
\begin{itemize}
\item Text measures depend on disclosure choices, which change over time and are endogenous.
\item Causality is not identified; firm FE mitigate but do not remove concern.
\item Climate-risk scores may capture generic ``risk'' rather than climate specifically.
\end{itemize}
\end{alertbox}

\section*{Round 2: Level 1 Fill-in}
\begin{enumerate}
\item Measurement: text in \blank{1cm}-K filings, supervised ML.
\item Three risk categories: physical, \blank{1.5cm}, regulatory.
\item Cash effect: +\blank{1cm}--\blank{1cm} pp of assets.
\item Capex effect: $-$\blank{1cm}--\blank{1cm} pp.
\item Sample years: \blank{1cm}--\blank{1cm}.
\end{enumerate}

\section*{Round 3: Level 2 Prompts}
\begin{enumerate}
\item Discuss the text-measurement methodology and validation strategy.
\item Why do firms increase cash and cut capex in response to climate risk?
\item Relate LSTY to BKS (2009) and FRS (2021).
\item Discuss causal identification in this kind of text-based climate-risk panel.
\end{enumerate}

\section*{Round 4: Minimal Scaffolding}
From a blank page: (i) measurement, (ii) validation, (iii) spec, (iv) headline magnitudes, (v) caveats.

\section*{Round 5: Blank Page}
Essay: how do firms respond to climate risk? LSTY evidence.

\vspace{6pt}
\begin{drillbox}
\textbf{Speed Drill.} (1) Text source. (2) Three risk categories. (3) Cash response. (4) Capex response. (5) Sample size.
\end{drillbox}

\section*{Quick Reference \& Answer Key}
\begin{answerbox}
\begin{itemize}
\item \textbf{Measure.} 10-K text, supervised ML, three risk subtypes.
\item \textbf{Cash.} Rises modestly with climate risk.
\item \textbf{Capex.} Falls.
\item \textbf{Leverage.} Declines modestly.
\end{itemize}
\end{answerbox}

\begin{keybox}
\textbf{30-second exam pitch.} Li, Shan, Tang \& Yao (2024) build text-based firm-level climate-risk scores from 10-K filings, split into physical, transition, and regulatory categories, validated against objective and disclosure-based exposures. In a US panel of $\sim$6{,}000 firms 2000--2020, firms with higher lagged climate-risk score hold more cash ($+0.5$--$1$ pp of assets), cut capex ($-0.3$--$0.5$ pp), and reduce leverage. Physical-risk-exposed firms drive most of the effect. The interpretation is that firms actively manage climate risk via financial flexibility (BKS 2009 precautionary cash extended to a new risk dimension) and via investment option value (wait-and-see under regulatory uncertainty). The paper is a leading template for text-based climate-risk measurement and shows that climate risk is already reshaping real corporate decisions.
\end{keybox}

\end{document}
